\section{Сквозные сценарии взаимодействия}

\subsection{Публикация текущих данных при активном пользователе}

\begin{minted}{text}
sequenceDiagram
  participant FE as Frontend
  participant FB as Firebase
  participant MCU as ESP8266

  FE->>FB: set /status/last_viewed = now (каждые 30с)
  MCU->>FB: get /status/last_viewed (каждые 10с)
  FB-->>MCU: last_viewed
  MCU->>MCU: diff = |epoch - last_viewed|
  alt diff < 120s
    MCU->>FB: set /current = {t,h,g}
    FB-->>FE: onValue(/current) update
  else diff >= 120s
    MCU->>FB: remove /current
    FB-->>FE: /current = null
  end
\end{minted}

\subsubsection*{Инженерный смысл сценария}

Сценарий реализует адаптивную передачу данных. Если пользователь недавно взаимодействовал с интерфейсом, система поддерживает активный live-поток. Если пользователь неактивен, поток отключается автоматически, что уменьшает число операций записи/чтения и экономит ресурсы сети.

\subsection{Почасовая запись истории}

\begin{minted}{text}
sequenceDiagram
  participant MCU as ESP8266
  participant FB as Firebase
  participant FE as Frontend

  MCU->>MCU: Каждый цикл проверка epoch % 3600 < 60
  MCU->>FB: set /history/{hourStartEpoch} = {t,h,g}
  FE->>FB: get /history (по кнопке Load data)
  FE->>FE: фильтрация по диапазону и сортировка
  FE->>FE: рендер графиков
\end{minted}

\subsubsection*{Инженерный смысл сценария}

Исторический поток отделён от live-потока и имеет другую частоту. Это предотвращает разрастание базы «мелкими» записями и сохраняет достаточную аналитическую ценность для графиков и отчётности.

\section{Машина состояний прошивки (упрощённо)}

\begin{minted}{text}
stateDiagram-v2
  [*] --> Boot
  Boot --> DisplayInit
  DisplayInit --> WifiTry
  WifiTry --> WifiOk: connected
  WifiTry --> Offline: timeout
  WifiOk --> FirebaseInit
  FirebaseInit --> SensorInit
  Offline --> SensorInit
  SensorInit --> Monitoring
  Monitoring --> Monitoring: read sensors + render OLED
  Monitoring --> CloudSync: Wi-Fi && FirebaseReady && validTime
  CloudSync --> Monitoring
\end{minted}

\subsection{Пояснение по состояниям}

\begin{itemize}
    \item \textbf{Boot / DisplayInit} --- аппаратная готовность и первичная верификация канала локальной индикации.
    \item \textbf{WifiTry / WifiOk / Offline} --- развилка сетевой доступности без блокировки основного назначения устройства.
    \item \textbf{FirebaseInit} --- подготовка облачного канала только при наличии сетевой инфраструктуры.
    \item \textbf{SensorInit / Monitoring} --- переход к целевому рабочему режиму наблюдения.
    \item \textbf{CloudSync} --- условный контур облачного обмена, активируемый только при валидных предусловиях.
\end{itemize}

\section{Важные инженерные особенности}

\begin{enumerate}
    \item \textbf{UTC как единая временная шкала}: NTP offset = \texttt{0}, timestamps хранятся в epoch-секундах.
    \item \textbf{Асинхронный Firebase-контур}: сетевые операции не блокируют цикл измерений и индикации.
    \item \textbf{Событийное управление live-потоком}: \texttt{/current} зависит от heartbeat-сигнала пользователя.
    \item \textbf{Экономичный исторический контур}: архив формируется батчево и запрашивается по требованию.
    \item \textbf{Явный fail-fast на OLED}: в текущей конфигурации отсутствие дисплея считается критическим отказом.
\end{enumerate}

\section{Ограничения текущей реализации}

\begin{enumerate}
    \item \textbf{Секреты в \texttt{secrets.h}}: учётные данные хранятся в коде и требуют усиления политики безопасности.
    \item \textbf{Чтение всей ветки истории}: при росте объёма данных возрастает нагрузка на сеть и клиентскую фильтрацию.
    \item \textbf{Отсутствие device-auth}: модель \texttt{NoAuth} снижает уровень защищённости подключения устройства.
    \item \textbf{Жёсткая реакция на сбой экрана}: отсутствует деградированный режим работы без OLED.
\end{enumerate}

\section{Рекомендации по развитию}

\begin{enumerate}
    \item Вынести секреты в приватный конфиг и исключить из репозитория.
    \item Добавить аутентификацию устройства (custom token через доверенный backend).
    \item Перейти к диапазонным запросам истории и индексированию по времени.
    \item Внедрить буферизацию оффлайн-измерений с последующей отправкой после восстановления сети.
    \item Добавить телеметрию состояния устройства: uptime, RSSI, ошибки публикации, статистику ретраев.
\end{enumerate}

\section{Карта соответствия «папка → ответственность»}

\begin{longtable}{|p{0.47\textwidth}|p{0.43\textwidth}|}
\hline
\textbf{Путь} & \textbf{Ответственность} \\
\hline
\path{course_work/src/main.cpp} & orchestration boot + runtime loop \\
\hline
\path{course_work/src/sensor_service.cpp} & чтение DHT11 и MQ-2, калибровка MQ-2 \\
\hline
\path{course_work/src/network_service.cpp} & Wi-Fi + NTP + валидация времени \\
\hline
\path{course_work/src/firebase_service.cpp} & async read/write Firebase, activity logic \\
\hline
\path{course_work/src/display_service.cpp} & визуализация статусов и телеметрии на OLED \\
\hline
\path{iot-dashboard/src/hooks/useCurrentData.ts} & realtime-подписка на \texttt{current} + heartbeat \\
\hline
\path{iot-dashboard/src/hooks/useHistoryData.ts} & загрузка и фильтрация \texttt{history} \\
\hline
\path{iot-dashboard/src/components/CurrentData.tsx} & карточки текущих показаний \\
\hline
\path{iot-dashboard/src/components/HistoryCharts.tsx} & графики истории и фильтры диапазона \\
\hline
\path{iot-dashboard/src/utils/telemetry.ts} & парсинг payload и нормализация газа \\
\hline
\end{longtable}

\section{Итоговое заключение}

Реализованная система образует завершённый технический контур IoT-мониторинга: устройство стабильно собирает телеметрию, облачный слой согласованно хранит и синхронизирует данные, а frontend предоставляет пользователю актуальную и историческую картину состояния среды.

Ключевое достоинство архитектуры --- предсказуемое поведение под нагрузкой за счёт разделения потоков данных, асинхронного обмена и строгой временной модели на базе UTC epoch.
